\section{Discussion}

The central finding is that for trace-based verdict prediction, learning efficiency strongly correlates with information density rather than information quantity.
Exception events concentrate the runtime's primary fault signal---error type, stack frame, and message---in a compact format (median 150 tokens), while non-exception components (variable states, loop iterations, control flow) record \emph{how} the program ran, not \emph{whether} it failed.

This density perspective explains three key results: (1)~adding trace information beyond exceptions provides no benefit (Table~\ref{tab:ablation}), because low-density content competes for attention; (2)~a tail-truncation budget sweep reveals a sharp phase transition at ${\sim}$500 tokens (Table~\ref{tab:tail_cross_model}), matching the boundary where diagnostic events dominate; (3)~the SFT advantage over regex concentrates on caught-exception boundary cases, confirming semantic signal extraction (\S\ref{sec:ablation}).
The narrow cross-model gap (Table~\ref{tab:cross_model}) supports that signal density, not model capacity, is the binding constraint.

\paragraph{Why exception\_only outperforms ae\_only.}
Although \texttt{ae\_only} is a strict superset, it underperforms \texttt{exception\_only} (Table~\ref{tab:ablation}), indicating extra content introduces noise rather than signal.

\paragraph{SFT vs.\ regex baselines.}
A natural question is whether SFT is necessary when exception signals are so discriminative.
We evaluate two regex baselines: detecting \texttt{[ERROR]} tokens ($B_{\text{error}}$, 93.45\%)
and detecting the \texttt{[No exception data in trace]} marker ($B_{\text{marker}}$, 99.33\%).
The optimal regex still misclassifies 8 samples, while SFT misclassifies only 2
(McNemar $p{=}0.014$). All 6 discordant samples involve \emph{caught exceptions}---programs
where a \texttt{try/except} block catches a \texttt{ValueError} but execution succeeds.
SFT learns to distinguish caught from uncaught exceptions; regex cannot.
Beyond accuracy, SFT provides calibrated probability estimates (ECE${=}$0.003),
format robustness (96.55\% without markers; Appendix~\ref{sec:nomarker}), and requires no task-specific engineering.

\paragraph{Why zero-shot and few-shot prompting fall short.}
Zero-shot prompting spans a 42pp range across models (\S\ref{sec:additional}), which SFT compresses to sub-1pp; few-shot ICL improves marginally, confirming the bottleneck is verdict calibration rather than exception recognition.


\paragraph{Near-perfect floor, not hard ceiling.}
The 2 errors (both \texttt{GeneratorExit}; Appendix~\ref{sec:error_overlap}) are shared across configurations; CodeLlama-13B seed 123 achieves 100\%, confirming 99.83\% is a convergence \emph{floor}, signal-limited rather than capacity-limited.

\paragraph{Tokenizer efficiency.}
CodeLlama's tokenizer produces ${\sim}$20\% more tokens than Qwen's on identical inputs, and a label-only loss variant collapses at tail-2048, indicating LoRA long-sequence instability (Appendix~\ref{sec:tokenizer_analysis}).

\paragraph{Tier hierarchy and methodological implications.}
The two-tier structure (Table~\ref{tab:cross_model}) suggests a fallback design: extract exception events first, resort to Tier~2 only when no exception is raised.
Loop benefit varies by architecture (Qwen-7B: +10.9pp; CodeLlama-13B: near-zero), likely reflecting differences in pre-training coverage of loop patterns.
The methodological insight---that information density strongly correlates with learning efficiency---transfers to any trace-based task where compact diagnostic events exist within verbose traces.

\paragraph{Density and cross-model robustness.}
Tier~1 and Tier~2 exhibit contrasting cross-model stability (Table~\ref{tab:cross_model}): \texttt{exception\_only}'s range is ${<}$1pp, while \texttt{loop\_only} spans ${\sim}$9pp.
However, $N{=}3$ models is insufficient to establish a continuous relationship, and Tier~1 signals are both high-density \emph{and} direct fault indicators, so density cannot be disentangled from signal type.
Applying label-only loss to DeepSeek full traces (91.7\% truncated) causes complete collapse (28.5\%), confirming that trace content density, not loss function, is the binding constraint.


\paragraph{Caveats on the density decomposition.}
Our 24\%/76\% gradient dilution vs.\ information overload decomposition
has two important caveats. First, the two factors are not orthogonal:
\texttt{label\_only} simultaneously changes the gradient source distribution
and the effective learning rate, so the decomposition is an approximation
rather than a causal isolation. Second, the decomposition is verified only
on the Qwen family (two scales); CodeLlama and DeepSeek \texttt{label\_only}
conditions collapse entirely, preventing cross-family replication.
Sequence length is an additional confounding variable that our design does not isolate.

\paragraph{Exception events capture test verdict signals, not just crashes.}
Although 55.5\% of failures are silent logic errors, \texttt{exception\_only} achieves 100\% fail recall because test assertions generate \texttt{AssertionError} events, confirming ``exception events'' encompasses both crashes and assertion failures.

\begin{figure}[t]
\centering
\includegraphics[width=\columnwidth]{training_dynamics}
\caption{Training dynamics comparison. Exception-only (blue) converges to 99.33\% within 50 steps and saturates at 99.83\% by step 150. Full-trace with label-only loss (red, dashed) oscillates between 82\% and 89\% across all 444 training steps without stable convergence, consistent with overfitting to spurious patterns in verbose trace content.}
\label{fig:training_dynamics}
\end{figure}

\paragraph{Information density and learning dynamics.}
Figure~\ref{fig:training_dynamics} shows the contrast: \texttt{exception\_only} converges rapidly, while \texttt{full\_trace} with label-only loss oscillates throughout---a signature of overfitting to spurious patterns (Appendix~\ref{sec:training_dynamics}).
Seed variance negatively correlates with input density (Table~\ref{tab:ablation}).
